%\ifodd\value{page}\hbox{}\newpage\fi
\chapter*{Śrī Lalitā Sahasranāma — Śloka 18}
\addcontentsline{toc}{chapter}{Śloka 18 - Nomi 36,37}

\begin{sloka}
{\sanskritfont
\begin{tabular}{c}
इन्द्रगोपपरिक्षिप्तस्मरतूणाभजङ्घिका ।\\
गूढगुल्फा कूर्मपृष्ठजयिष्णुप्रपदान्विता ॥१८॥
\end{tabular}

\vspace{1.5em}

{\middlesize \iastfont
indragopa-parikṣipta-smara-tūṇābha-jaṅghikā |\\
gūḍha-gulphā kūrma-pṛṣṭha-jayiṣṇu-prapadānvitā ||18||}

\vspace{1em}

{\middlesize
\anvaya{%
इन्द्रगोपपरिक्षिप्तस्मरतूणाभजङ्घिका ।
गूढगुल्फा कूर्मपृष्ठजयिष्णुप्रपदान्विता ।
}
}

\middlesize \iastfont \itmean{%
«Le cui gambe sono simili alla faretra di \textit{Smara}, circondate dal rosso degli \textit{indragopa};  
dalle caviglie ben raccolte, dotata di piedi che superano in grazia il dorso della tartaruga.»%
}

}
\end{sloka}

\wordmeaning{%
\wordline{इन्द्रगोप-परिक्षिप्त-स्मर-तूण-आभ-जङ्घिका}
{indragopa-parikṣipta-smara-\\tūṇābha-jaṅghikā}
{ le cui gambe (\textit{jaṅghikā}) sono simili (\textit{ābha}) alla faretra (\textit{tūṇa}) di \textit{Smara}, circondate (\textit{parikṣipta}) dal rosso vivo degli \textit{indragopa}, piccoli insetti vermigli frequentemente usati come termine di paragone poetico; }%

\wordline{गूढ-गुल्फा कूर्म-पृष्ठ-जयिष्णु-प्रपद-अन्विता}
{gūḍha-gulphā kūrma-pṛṣṭha-\\jayiṣṇu-prapadānvitā}
{ dalle caviglie (\textit{gulpha}) ben raccolte e non prominenti (\textit{gūḍha}), dotata (\textit{anvitā}) di piedi (\textit{prapada}) che superano (\textit{jayiṣṇu}) il dorso (\textit{pṛṣṭha}) della tartaruga (\textit{kūrma}) in grazia e bellezza; }%
}

\textit{Śloka} 18 prosegue la contemplazione discendente soffermandosi sulle gambe, le caviglie e i piedi. Le gambe sono paragonate alla faretra di \textit{Smara/Kāma}, circondata dal rosso degli \textit{indragopa}: un’immagine che unisce colore, tensione e potenza latente, senza alludere a un gesto esplicito. Il desiderio è evocato come energia trattenuta, non come azione.

Nel secondo \textit{pāda}, l’attenzione si sposta sulle caviglie e sui piedi. Le caviglie sono dette “nascoste”, cioè raccolte e compatte, prive di asperità. I piedi sono descritti come capaci di superare persino il dorso della tartaruga, tradizionale termine di paragone per stabilità e grazia. Il confronto con il dorso della tartaruga può anche alludere a \textit{Kūrma}, figura del sostegno cosmico: i piedi della Dea sono così intesi come fondamento stabile, ciò che regge e rende possibile il manifestarsi del mondo.

Nel linguaggio di questo \textit{śloka}, \textit{Smara} e \textit{Kūrma} non introducono una mitologia narrativa, ma indicano due principi complementari. \textit{Smara} non è il desiderio come impulso erotico consumato, ma come tensione vitale che risveglia e mette in movimento: la faretra richiama una potenza presente e trattenuta, pronta ma non ancora scagliata. \textit{Kūrma}, al contrario, è figura del sostegno e della stabilità: ciò che regge, porta e rende possibile il manifestarsi senza crollo. Collocati rispettivamente alle gambe e ai piedi, questi due riferimenti delineano una dinamica essenziale: il desiderio che muove e il fondamento che sostiene. 

Questo \textit{śloka} contiene due \textit{nāma}.  
\textbf{\textit{Indragopa-parikṣipta-smara-tūṇābha-jaṅghikā}} presenta le gambe come luogo di potenza trattenuta e colore vitale.  
\textbf{\textit{Gūḍha-gulphā kūrma-pṛṣṭha-jayiṣṇu-prapadānvitā}} raffigura caviglie e piedi come segni di compattezza, stabilità e grazia superiore.
